\scanpage{91}
% archive foliation: 90
La donnée d'un tel \deleted{(indép. de $i$)} $c_\Omega$ équivaut à une
\unsure{restriction|réduction} du Groupe structural de
$\uline{\mathrm{Conn}}_i\unsure{\xi|\varepsilon}$ à $\Omega\otimes\uline\omega$, ou encore
(indépendant de $i$) la donnée d'un $c_\Omega$ pour lequel
$(\unsure{A|*})$ est lin et $(**)$ est \unsure{vérif.|vérifié}, équivaut
à une restriction du Groupe structural
$\mathfrak g\otimes\omega$ de \uline{Conn} à $\Omega\otimes\uline\omega$.

On s'intéresse au cas où \deleted{\unsure{si|C}} $c_\Omega$ est
tel que la composée
\diagram{\begin{tikzcd}
(-2i):\Omega\arrow[r]&\mathcal L\otimes\uline\omega\arrow[r,dash]&\uline\omega
\end{tikzcd}}
\inserted{(\uline{ce qui implique $2$ inv.})}
soit un \uline{iso}. Alors le triplet
$(\mathfrak g,\mathcal L,c_\Omega)$
(\unsure{où|ou} ce qui revient au
$=$, $(\mathcal L\supset\Omega,c_\Omega)$) est défini à isom
unique près par $X/S$ (et existe ssi
$\uline\omega=\underline\Omega^1_{X/S}$ est inversible) comme
\[
\mathcal L\simeq P^1_{X/S}(\uline\omega)\otimes\uline\omega^{-1}
\qquad(\text{d'où }\mathcal L\otimes\uline\omega\simeq P^1_{X/S}(\uline\omega))
\]
et
\diagram{\begin{tikzcd}
c_\Omega:\Omega\simeq\uline\omega\arrow[r]&\mathcal L\otimes\uline\omega\simeq P^1_{X/S}(\uline\omega)
\end{tikzcd}}
n'est autre que
$d^1_{X/S,\underset{\uline\omega}{\unsure{\omega|\infty}}}$ --

\figure{\begin{tikzpicture}[x=1cm,y=1cm]
\draw (0,7)--(0,2.2) .. controls (0,.9) and (.2,.4) .. (1.4,-.1)
.. controls (2.4,-.5) and (2.9,-1.1) .. (4.5,-1.5);
\draw (2.1,4.4) .. controls (1.9,4.9) and (2.6,4.9) .. (3.4,4.9)
--(5.4,5.1);
\end{tikzpicture}}

